% rtwaterflow — Benutzerhandbuch (LaTeX-Quelle für das gesetzte PDF)
% Bauen:  latexmk -pdf Benutzerhandbuch.tex     (oder: pdflatex zweimal)
% Der in der Anwendung unter /manual eingebettete Kurztext ist Benutzerhandbuch.md;
% dieses Dokument ist das ausführliche, gesetzte Handbuch mit derselben Tiefe wie
% das Schwesterprojekt rtpowerflow.
\documentclass[11pt,a4paper,parskip=half]{scrreprt}

\usepackage[T1]{fontenc}
\usepackage[utf8]{inputenc}
\usepackage{lmodern}
\usepackage[ngerman]{babel}
\usepackage{microtype}
\usepackage{amsmath}
\usepackage{booktabs}
\usepackage{longtable}
\usepackage{tabularx}
\usepackage{xcolor}
\usepackage{enumitem}
\usepackage{listings}
\usepackage{graphicx}
\graphicspath{{img/}}
\usepackage{tikz}
\usetikzlibrary{positioning,arrows.meta,calc,shapes.geometric,fit,backgrounds}
\usepackage[most]{tcolorbox}
\usepackage[colorlinks=true,linkcolor=blue!50!black,urlcolor=blue!50!black,
            citecolor=blue!50!black]{hyperref}

% --- Farben (Wasser-Blau als Leitfarbe) ---
\definecolor{wasser}{RGB}{0,105,148}
\definecolor{wasserhell}{RGB}{224,240,247}
\definecolor{shell}{gray}{0.96}
\definecolor{gruen}{RGB}{60,140,60}
\definecolor{amber}{RGB}{200,140,0}
\definecolor{rot}{RGB}{190,40,40}

% --- Eigene Befehle ---
\newcommand{\ui}[1]{\textsf{\bfseries #1}}      % Oberflächen-Beschriftung
\newcommand{\api}[1]{\texttt{#1}}               % REST-/WS-Endpunkt
\newcommand{\file}[1]{\texttt{#1}}              % Datei / Pfad

% --- Beispiel-/Merkkasten (blau umrandet, umbruchfähig) ---
\newtcolorbox{beispiel}[1][Beispiel]{%
  breakable, enhanced, colback=wasserhell, colframe=wasser,
  boxrule=0pt, leftrule=3pt, arc=1pt, left=8pt, right=8pt, top=5pt, bottom=5pt,
  fonttitle=\bfseries\sffamily, coltitle=wasser, title={#1}}

% --- Code-Listings ---
\lstset{
  basicstyle=\ttfamily\small, backgroundcolor=\color{shell},
  frame=single, framesep=4pt, rulecolor=\color{gray!40},
  breaklines=true, columns=fullflexible, keepspaces=true,
  showstringspaces=false, extendedchars=true, inputencoding=utf8,
  literate={ä}{{\"a}}1 {ö}{{\"o}}1 {ü}{{\"u}}1 {Ä}{{\"A}}1 {Ö}{{\"O}}1
           {Ü}{{\"U}}1 {ß}{{\ss}}1 {→}{{$\rightarrow$}}1 {≥}{{$\geq$}}1
           {≤}{{$\leq$}}1 {·}{{$\cdot$}}1 {°}{{$^\circ$}}1}

\emergencystretch=3em
\hyphenation{Trink-wasser-netz Druck-zonen Hoch-be-hälter Nach-frage-Engine
             Zu-stands-schätzung Schlecht-punkt}

\title{rtwaterflow\\[0.3em]\Large Benutzerhandbuch}
\subtitle{Echtzeit-Simulation von Trinkwasser-Verteilnetzen\\für Lehre und
          Demonstration}
\author{}
\date{Stand: Juli 2026}

\begin{document}
\maketitle

\begin{tcolorbox}[colback=gray!5,colframe=gray!50,arc=1pt,boxrule=0.5pt]
\small
\textbf{KI-erzeugte Software.} Quelltext, Tests und Dokumentation dieser
Plattform — einschließlich dieses Handbuchs — wurden von einem KI-Programmier\-agenten
(Claude Code, Anthropic) unter menschlicher Leitung geschrieben: eine Person hat
Anforderungen und Fachentscheidungen vorgegeben, die Ergebnisse geprüft und
freigegeben. rtwaterflow ist ein Forschungs- und Lehrprojekt; die mitgelieferten
Netze sind synthetisch bzw.\ aus offenen Geodaten synthetisiert und modellieren
kein konkretes reales Versorgungsgebiet. \textbf{Lizenz:} MIT (Quelltext und
Dokumentation); die Hydraulik-Engine \href{https://github.com/e2nIEE/pandapipes}{pandapipes}
ist eine permissiv lizenzierte (BSD-3) Abhängigkeit.
\end{tcolorbox}

\tableofcontents

% ===================================================================
\chapter{Einführung}
\label{ch:einfuehrung}

\section{Was ist rtwaterflow?}

rtwaterflow ist eine Lehr- und Forschungsplattform, die ein kommunales
\textbf{Trinkwasser-Verteilnetz in beschleunigter Echtzeit} simuliert. Auf einer
Karte werden Drücke, Fließgeschwindigkeiten und Fließrichtungen sichtbar;
fallende Behälterstände, Druckabfälle an Hausanschlüssen, trockenfallende Brunnen
und die DVGW-Regelverstöße in der Alarmzentrale sind die zentralen Lehreffekte.

Die Plattform ist der Kaltwasser-Zwilling von
\href{https://github.com/markisbell/rtheatflow}{rtheatflow} (Fernwärme) und
teilt sich mit diesem und dem Ausgangsprojekt
\href{https://github.com/markisbell/rtpowerflow}{rtpowerflow} (Strom) einen
gemeinsamen Plattformkern: eine \emph{einmal bauen, günstig rechnen}-Engine, ein
dreischichtiges Beobachtbarkeitsmodell, eine WebSocket-Datenleitung, eine
Aufzeichnung/Export-Schicht und eine React/Leaflet-Oberfläche. Die Unterschiede
liegen allein in der Physik-Domäne — hier: der stationären Rohrnetzhydraulik nach
\textbf{Darcy-Weisbach}, gerechnet mit
\href{https://github.com/e2nIEE/pandapipes}{pandapipes~0.14.0}.

\section{Die drei Anwendungen}

\begin{description}[style=nextline,leftmargin=1.2em]
\item[Live-Simulation] Das Herzstück: ein Netz läuft Takt für Takt, jeder Takt
  ist eine vollständige, stationäre Hydraulikrechnung. Sie beobachten das Netz in
  drei Sichten (Realität, Gemessen, Schätzung), lösen Ereignisse aus (Löschwasser,
  Rohrbruch, Leckage) und stellen Betriebsmittel (Pumpen, Behälter) und Umwelt
  (Hitzetag, Dürre) ein.
\item[NetzStudio-Editor] Zeichnen Sie ein eigenes Netz auf \emph{echten}
  OpenStreetMap-Straßen, prüfen Sie es live gegen die drei DVGW-Lastfälle (W~400-1)
  und übernehmen Sie es in die Bibliothek (Kapitel~\ref{ch:editor}).
\item[Geodaten-Bündelbauer] Ein Offline-Werkzeug, das eine echte deutsche Stadt
  aus OSM-Straßen und einem Höhenmodell in ein Netzbündel verwandelt
  (Kapitel~\ref{ch:geodaten}). Drei mitgelieferte Netze (\file{alpen},
  \file{neubeuern}, \file{kevelaer}) sind so entstanden.
\end{description}

\section{Das Zeitmodell}

Ein Simulationsschritt entspricht \textbf{einer Minute Netzzeit}. Die Engine
rechnet eine feste Zahl Schritte je Tag (Standard 1440 = ein Schritt je
Minute) und wickelt am Tagesende auf einen neuen Tag um — so laufen die
Tages-, Wochen-, Saison- und Wetterprofile der Nachfrage zyklisch durch. Die
Wandzeit je Schritt ist voreingestellt auf eine reale Sekunde je Netzminute (ein
Simulationstag in rund 24 Minuten); die Hydraulikrechnung selbst braucht nur
einen Bruchteil dieser Sekunde, sodass deutlich schnellere Läufe möglich sind.

\begin{center}
\begin{tabular}{@{}lll@{}}
\toprule
\ui{Schritte/Tag} & \ui{Wandzeit/Schritt} & Charakter \\
\midrule
1440 & 1,0 s & ein Netztag in 24 Minuten (Standard) \\
96   & 1,0 s & 15-Minuten-Raster, schneller Durchlauf \\
1440 & 0,05 s & Zeitraffer für lange Verläufe \\
\bottomrule
\end{tabular}
\end{center}

Jeder Takt ist eine \textbf{stationäre Momentaufnahme}. Das Modell ist damit
quasistatisch: Druckstöße/Wasserschläge (DVGW~W~303) werden bewusst nicht
gerechnet (siehe Kapitel~\ref{ch:grenzen}).

\section{Validierte Physik}
\label{sec:validierung-intro}

Ein Lernender muss darauf vertrauen können, dass die Zahlen auf der Karte echte
Hydraulik sind. Die Engine ist deshalb gegen das unabhängige Referenzwerkzeug
\textbf{EPANET} (über WNTR) kreuzvalidiert: die Standardnetze \textbf{Net1}
(9~Knoten) und \textbf{Net3} (92~Knoten), in pandapipes mit \api{swamee-jain}
nachgebaut, stimmen mit WNTRs EpanetSimulator auf $<0{,}001$~bar (Druck) bzw.\
$\sim0{,}5\,\%$ (Durchfluss) überein. Behälter- und Hydrantendynamik sowie die
druckabhängige Nachfrage sind gegen eigene EPANET/WNTR-Orakel geprüft, die
Nachfrage-Engine gegen den DVGW-Korridor W~410. Die Methodik steht ausführlich in
Anhang~\ref{app:validierung} und in \file{docs/BENCHMARKS.md}.

\section{Weiterführende Dokumentation}

\begin{description}[leftmargin=1.2em]
\item[\file{docs/ARCHITECTURE.md}] Die vollständige Systemarchitektur (Module,
  Datenfluss, die drei Schichten, Solver-Leiter, Frontend).
\item[\file{docs/BENCHMARKS.md}] Die EPANET/WNTR-Validierung mit Reproduktion.
\item[\file{docs/GEODATA\_BUILDER.md}] Der Geodaten-Bündelbauer im Detail.
\item[\file{docs/COMPLIANCE.md}] Jede DVGW-Regel der Alarmzentrale mit Zitatstelle.
\item[\file{docs/API.md}] Die vollständige REST-/WebSocket-Referenz (65 Routen).
\end{description}

% ===================================================================
\chapter{Installation und Start}
\label{ch:installation}

\section{Windows-Schnellstart (empfohlen)}

\file{start\_rtwaterflow.bat} startet Backend (Port~8002) und Oberfläche
(Port~5175) in zwei Konsolen und öffnet den Browser. \file{stop\_rtwaterflow.bat}
beendet alles wieder. Die Launcher erkennen ihre eigenen Prozesse anhand eines
Fenstertitels und einer Kommandozeilen-Markierung, damit sie fremde
Python-Prozesse nicht versehentlich beenden.

\section{Docker Compose}

\begin{lstlisting}[language=bash]
docker compose up -d
\end{lstlisting}

startet fünf Dienste: Backend (8002), Oberfläche (8082), InfluxDB (8088),
den Telemetrie-\emph{collector} und Grafana (3002).

\begin{center}
\begin{tabular}{@{}lll@{}}
\toprule
Dienst & Port & Zweck \\
\midrule
Backend (FastAPI) & 8002 & REST + WebSocket, die Simulations-Engine \\
Oberfläche (React) & 8082 & die Bedienoberfläche \\
InfluxDB & 8088 & Zeitreihen-Speicher \\
Grafana & 3002 & Langzeit-Dashboards (Kapitel~\ref{ch:grafana}) \\
\bottomrule
\end{tabular}
\end{center}

\begin{beispiel}[Portschema der Geschwisterprojekte]
netzsim 8000/5173 $\cdot$ rtheatflow 8001/5174 $\cdot$ \textbf{rtwaterflow
8002/5175} (Compose-Hostports 8002/8082/8088/3002). Mehrere Geschwister laufen so
kollisionsfrei nebeneinander.
\end{beispiel}

\section{Lokaler Start (Entwicklung)}

\begin{lstlisting}[language=bash]
python -m venv .venv
.venv\Scripts\pip install -r requirements.txt -e .[dev]
.venv\Scripts\python -m rtwaterflow.main        # Backend :8002
cd ui && npm install && npm run dev             # Oberfläche :5175
\end{lstlisting}

Die Oberfläche leitet \api{/api} und \api{/ws} auf \texttt{127.0.0.1:8002}
weiter. Der Extra \texttt{[dev]} zieht die Test-Werkzeuge (u.\,a.\ \texttt{wntr}
für die EPANET-Validierung); für den reinen Betrieb genügt
\file{requirements.txt}.

% ===================================================================
\chapter{Das Netz und die Netzbibliothek}
\label{ch:netz}

\section{Die fünf Dateien eines Netzbündels}

Ein Netz ist ein Bündel aus fünf JSON-Dateien im Verzeichnis
\file{data/networks/<id>/}. Die Formate sind pandapipes-nah und werden beim Laden
streng geprüft (Anhang~\ref{app:dateiformat}):

\begin{description}[leftmargin=1.2em]
\item[\file{network\_structure.json}] Ein Eintrag je Knoten mit \textbf{Geländehöhe}
  \texttt{elevation\_m} — der hydrostatischen Kernphysik: rund $0{,}098$~bar je
  Höhenmeter. Optional eine Quellenangabe (\texttt{attribution}) für Geodaten-Netze.
\item[\file{pipes.json}] Rohre mit Innendurchmesser und integraler Rauheit~$k$.
  Der Rohrkatalog leitet $k$ aus dem Material (PE/PVC 0{,}1; GGG/St/AZ 0{,}4;
  GG 1{,}0~mm, GW~303-1) und die lichte Weite aus DN bzw.\ der PE-d-Reihe ab.
\item[\file{consumers.json}] Verbraucher mit Archetyp und Größe; daraus baut die
  Nachfrage-Engine (Kapitel~\ref{ch:nachfrage}) realistische Profile. Bei
  Unterversorgung entnehmen sie druckabhängig weniger (PDA,
  Kapitel~\ref{ch:ereignisse}).
\item[\file{supply.json}] Einspeisungen (Festdruckknoten/Hochbehälter), Behälter
  mit Füllstandsdynamik, Pumpwerke, Druckminderer und Brunnenfelder
  (Kapitel~\ref{ch:behaelter}, \ref{ch:brunnen}).
\item[\file{environment.json}] Zeithorizont, Lufttemperatur, Tagestypen,
  Trockenheit und Jahreszeit für die Nachfrage- und Wettereffekte.
\end{description}

\section{Die mitgelieferte Bibliothek}
\label{sec:bibliothek}

Sieben Lehrnetze stehen bereit — vier synthetische Layouts und drei aus echten
Geodaten synthetisierte Netze. Kapitel~\ref{ch:rundgaenge} führt durch jedes
einzelne.

\begin{center}
\small
\begin{tabularx}{\linewidth}{@{}llrX@{}}
\toprule
Netz & Charakter & Knoten & Lehrpunkt \\
\midrule
Tutorial Hanglage & Tutorial & 5   & Hydrostatik; der Schlechtpunkt liegt am \emph{höchsten} Abnehmer \\
Musterdorf        & ländlich & 35  & zwei \textbf{Druckzonen} über einen Druckminderer; Pumpe + Hochbehälter; Hitzetag \\
Mustertal         & ländlich & 8   & \textbf{Gegenbehälter} — ein Abschnitt kehrt über den Tag die Richtung um \\
Lauenau           & ländlich & 8   & \textbf{Brunnen \& Grundwasser}; die Dürrekaskade lässt Haushalte trockenfallen \\
Alpen             & real     & 182 & echtes OSM/EU-DEM, flacher Niederrhein, eine Schwerkraftzone \\
Neubeuern         & real     & 215 & echtes OSM/EU-DEM, hügeliges Inntal, \textbf{PRV-Druckzonen} über $\sim78$~m Relief \\
Kevelaer          & real     & 544 & Stadtnetz — der Leistungs-Stresstest ($<100$~ms/Takt) \\
\bottomrule
\end{tabularx}
\end{center}

\begin{beispiel}[Ein eigenes Netz importieren]
Über \ui{NetzStudio} $\to$ \ui{Import} laden Sie ein eigenes Fünf-Datei-Bündel;
es erscheint in der Bibliothek mit dem Charakter \emph{user}. Alternativ zeichnen
Sie eines im Editor (Kapitel~\ref{ch:editor}). Der Loader prüft dabei denselben
Vertrag wie für die mitgelieferten Netze (Anhang~\ref{app:dateiformat}).
\end{beispiel}

% ===================================================================
\chapter{Arbeitsablauf: Netz — Nachfrage — Live}
\label{ch:arbeitsablauf}

Der typische Ablauf hat drei Schritte: ein Netz wählen, die Nachfrage/Umwelt
einstellen, die Live-Ansicht beobachten und steuern.

\begin{figure}[t]
\centering
\begin{tikzpicture}[
  node distance=6mm and 12mm, font=\small,
  box/.style={draw=wasser, thick, rounded corners=2pt, fill=wasserhell,
              minimum width=26mm, minimum height=9mm, align=center},
  arr/.style={-{Latex[length=2mm]}, thick, draw=wasser}]
\node[box] (bundle) {Fünf-Datei-\\Bündel (JSON)};
\node[box, right=of bundle] (loader) {\texttt{data\_loader}\\(prüfen)};
\node[box, right=of loader] (builder) {\texttt{network\_}\\\texttt{builder}};
\node[box, right=of builder] (sim) {\texttt{Simulator}\\1 min/Takt};
\node[box, below=10mm of sim] (store) {\texttt{StateStore}\\(Verlauf)};
\node[box, left=22mm of store] (ui) {Oberfläche\\(WebSocket)};
\draw[arr] (bundle) -- (loader);
\draw[arr] (loader) -- (builder);
\draw[arr] (builder) -- node[above,font=\scriptsize]{pandapipes-Netz} (sim);
\draw[arr] (sim) -- (store);
\draw[arr] (store) -- node[above,font=\scriptsize]{Frames} (ui);
\end{tikzpicture}
\caption{Der Datenpfad: einmal bauen, günstig rechnen. Jeder Takt überschreibt
nur die Einspeise-Spalten und löst neu — nie ein Neuaufbau.}
\label{fig:datenpfad}
\end{figure}

\section{Schritt 1: Netz wählen oder importieren}

In \ui{NetzStudio} wählen Sie ein Bündel aus der Bibliothek
(Abschnitt~\ref{sec:bibliothek}) oder importieren ein eigenes. Das gewählte Netz
wird \textbf{einmal} zu einem pandapipes-Netz gebaut; jeder Takt überschreibt
danach nur Injektionsspalten (Verbrauchs-Massenströme, Behälterhöhen,
Pumpenzustände) und rechnet warm gestartet neu (Abbildung~\ref{fig:datenpfad}).

\section{Schritt 2: Nachfrage und Umwelt}

Die Verbraucher tragen Archetyp und Größe; daraus erzeugt die Nachfrage-Engine
Tages-, Wochen- und Saisonprofile (Kapitel~\ref{ch:nachfrage}). Im Bereich
\ui{Umwelt} schalten Sie zwischen \ui{Normal} und \ui{Hitzetag} um — letzterer
hebt Temperatur und Trockenheit, verschiebt die Verbrauchsspitze in den Abend und
hebt die Nachtflüsse.

\section{Schritt 3: Die Live-Ansicht}

\subsection{Darstellungen}
Die Karte (Leaflet/OpenStreetMap) zeigt Knoten und Rohre georeferenziert. Die
\ui{Drucklinie} zeichnet zusätzlich das Höhen- und Druckprofil entlang eines
Pfades (Kapitel~\ref{ch:karte}). Tabellen listen Verbraucher, Behälter, Brunnen
und Alarme.

\subsection{Steuerung}
\ui{Start}/\ui{Pause} steuern den Taktlauf; einzelne Betriebsmittel schalten Sie
über ihre Regler (Pumpenmodi \ui{Auto}/\ui{Ein}/\ui{Aus}, Dürre-Schieber,
PDA-Umschalter). Eingaben werden optimistisch angezeigt, sodass auch im
pausierten Zustand Ihre Bedienung bestätigt wird.

\subsection{Die Übersicht lesen}
Die \ui{Übersicht} fasst die Schlüsselgrößen jedes Takts zusammen:

\begin{center}
\small
\begin{tabularx}{\linewidth}{@{}lX@{}}
\toprule
Größe & Bedeutung \\
\midrule
Schlechtpunkt & der niedrigste Versorgungsdruck über die \emph{Verbraucher}-Knoten \\
Massenbilanz & Einspeisung $-$ geliefert $-$ gespeichert $-$ Überlauf $-$ Export $-$ Entnahme \\
Fehlmenge & unbefriedigte Nachfrage (PDA) \\
Solver-Status & \texttt{ok} / \texttt{degraded} / \texttt{failed} \\
\bottomrule
\end{tabularx}
\end{center}

Ein niedriger Einspeisedruck (etwa der 0{,}5-bar-Auslauf eines Hochbehälters) ist
\emph{kein} Netzfehler — der Schlechtpunkt wird deshalb nur über Verbraucher
gebildet, nie über die Quelle.

\subsection{Interaktion mit Netzelementen}
Ein Klick auf ein Element öffnet ein Kontextmenü: einen Knoten anpinnen, einen
\ui{Wasserzähler} oder \ui{Drucksensor} setzen (Kapitel~\ref{ch:beobachtbarkeit}),
einen Verbraucher hinzufügen. Popups zeigen die ehrlichen Werte — z.\,B.\ Soll-
gegen Ist-Entnahme eines Verbrauchers, oder In$\to$Out-Druck eines
Druckminderers.

% ===================================================================
\chapter{Die Karte und die Ebenen}
\label{ch:karte}

\section{Druck- und Geschwindigkeitsebene}

Die Ebene \ui{Druck} färbt Knoten und Verbraucher nach dem lokalen Druck, an den
DVGW-Werten ausgerichtet:

\begin{center}
\begin{tabular}{@{}ll@{}}
\toprule
Farbe & Bedeutung \\
\midrule
\textcolor{rot}{\textbf{rot}}   & unter 2{,}0~bar (Mindestversorgungsdruck W~400-1) \\
\textcolor{gruen}{\textbf{grün}} & im Sollband 4--6~bar \\
\textcolor{amber}{\textbf{gelb}} & ab 8~bar (Ruhedruck) \\
\bottomrule
\end{tabular}
\end{center}

Die Ebene \ui{Geschwindigkeit} färbt Rohre (Warnung ab 2{,}0~m/s).
\textbf{Fließrichtungspfeile} zeigen den Durchfluss; in der gemessenen Sicht sind
sie ausgeblendet, wo keine Messung die Richtung belegt.

\section{Marker}
Eigene Markierungen tragen Einspeisungen, Druckminderer (mit Soll-/Ist-Druck),
Pumpwerke (mit Lauflampe und Rückschlag-Alarm), Hochbehälter (mit Füllstand und
Löschreserve-Band) und Hydranten/Rohrbrüche. Geodaten-Netze zeigen den OSM-/
EU-DEM-Quellennachweis in der Kartenecke.

\section{Die Drucklinie (hydraulische Höhenlinie)}

Die \ui{Drucklinie} zeichnet entlang eines Pfades die Geländehöhe und darüber die
Drucklinie (Energiehöhe) — die Stufe eines Druckminderers wird als sichtbarer
Sprung dargestellt. Der Pfad wird clientseitig per Dijkstra bestimmt (die Rohre
plus Pseudo-Kanten für Druckminderer und Pumpen).

\begin{figure}[t]
\centering
\begin{tikzpicture}[font=\small,scale=1.0]
% Achsen
\draw[->] (0,0) -- (10.3,0) node[right,font=\scriptsize]{Weg};
\draw[->] (0,0) -- (0,4.3) node[above,font=\scriptsize]{Höhe / m};
% Gelände
\draw[thick,brown!70!black] (0,3.4) -- (2,3.0) -- (4,2.3) -- (5.0,2.1)
      -- (7,1.4) -- (9.5,0.7);
\node[brown!70!black,font=\scriptsize] at (8.8,1.15) {Gelände};
% Drucklinie oben (Hochzone)
\draw[thick,wasser] (0,4.0) -- (2,3.85) -- (4,3.55) -- (5.0,3.45);
% PRV-Sprung
\draw[dashed,gray] (5.0,3.45) -- (5.0,2.75);
\node[font=\scriptsize,align=center] at (5.9,3.1) {Druck-\\minderer};
% Drucklinie unten (Tiefzone)
\draw[thick,wasser] (5.0,2.75) -- (7,2.4) -- (9.5,1.9);
\node[wasser,font=\scriptsize] at (2.0,4.15) {Drucklinie};
% Druck = Differenz
\draw[<->,gray] (7,1.4) -- (7,2.4) node[midway,right,font=\scriptsize]{Druck};
% Quelle
\filldraw[wasser] (0,4.0) circle (2pt) node[above left,font=\scriptsize]{Hochbehälter};
\end{tikzpicture}
\caption{Die Drucklinie: der senkrechte Abstand zwischen Drucklinie und Gelände
ist der Versorgungsdruck. Der Druckminderer setzt die Drucklinie der Tiefzone
herab und schützt sie vor Überdruck.}
\label{fig:drucklinie}
\end{figure}

% ===================================================================
\chapter{Beobachtbarkeit und Zustandsschätzung}
\label{ch:beobachtbarkeit}

Der pädagogische Kern der Plattform ist die \textbf{Beobachtbarkeit}: der
Unterschied zwischen dem, was physikalisch \emph{wahr} ist, dem, was ein Betreiber
tatsächlich \emph{misst}, und dem, was sich daraus \emph{schätzen} lässt.

\section{Die drei Sichten}

\begin{description}[leftmargin=1.2em]
\item[Realität] die vollständige Physik — jeder Knotendruck, jede Geschwindigkeit,
  jede verborgene Leckage.
\item[Gemessen] nur was Wasserzähler (Durchfluss und Druck), Drucksensoren und die
  Leitwarte wirklich liefern — mit Live- gegen 15-Minuten-Treue und ehrlichem
  Kaltstart (ein neuer Zähler zeigt bis zum ersten Fensterabschluss nichts an).
\item[Schätzung] der \textbf{Digital-Twin-Beobachter}: ein \emph{zweites} Netz,
  allein aus Operatorwissen rekonstruiert.
\end{description}

\begin{figure}[t]
\centering
\begin{tikzpicture}[font=\small, node distance=8mm,
  v/.style={draw, thick, rounded corners=2pt, minimum width=30mm,
            minimum height=13mm, align=center}]
\node[v, fill=gray!8, draw=gray!60] (r) {\textbf{Realität}\\alle Drücke\\+ verborgene Leckage};
\node[v, right=12mm of r, fill=wasserhell, draw=wasser] (m) {\textbf{Gemessen}\\nur Zähler +\\Sensoren + SCADA};
\node[v, right=12mm of m, fill=gruen!12, draw=gruen] (s) {\textbf{Schätzung}\\zweites Netz aus\\SCADA + Prioris};
\draw[-{Latex[length=2mm]}, thick, gray] (r) -- node[above,font=\scriptsize]{Filter} (m);
\draw[-{Latex[length=2mm]}, thick, wasser] (m) -- node[above,font=\scriptsize]{Vorwärtslösung} (s);
\end{tikzpicture}
\caption{Von der Realität bleibt in der gemessenen Sicht nur, was Sensoren
liefern; die Schätzung rekonstruiert daraus ein eigenständiges Netz.}
\label{fig:sichten}
\end{figure}

\section{Messgeräte platzieren}

Über das Kontextmenü setzen Sie \ui{Wasserzähler} an Verbrauchern (liefern
Durchfluss \texttt{mdot} und lokalen Druck \texttt{p}) und \ui{Drucksensoren} an
Knoten. Presets belegen typische Konstellationen — z.\,B.\ das \ui{scada}-Preset:
kritische Druck-Logger an Quelle und Netzenden, ohne Hausanschlusszähler. Die
Einspeisung ist als Leitwarten-Telemetrie stets gemessen.

\section{Der Beobachter (Digital Twin)}

pandapipes besitzt keinen Zustandsschätzer. Die geschätzte Sicht ist deshalb kein
Filter über der Wahrheit, sondern eine \textbf{unabhängige Vorwärtslösung}, die
nur mit Operatorwissen gespeist wird: Quell-/Behälter-/Pumpen-/Druckminderer-Tele\-metrie
(aus der laufenden Fahrweise übernommen) und die \emph{erwartete}, rauschfreie
Nachfrage für ungezählte Verbraucher. Jede Emitter-Entnahme (Leckage, Bruch,
Hydrant) wird im Zwilling zu Null gesetzt.

\begin{beispiel}[Die Ehrlichkeit des Beobachters]
Ein Rohrbruch an einem \emph{ungezählten} Knoten bewegt die gesamte Schätzung um
$\approx0$ — er bleibt unsichtbar. Dieselbe Anomalie an einem \emph{gezählten}
Verbraucher schlägt durch, weil der Zähler die Mehrentnahme liefert. Genau diese
Asymmetrie ist die Lehre zur Beobachtbarkeit; sie wird von Prüf-Tests
(\emph{tripwires}) abgesichert. Die Abweichung zwischen Zwilling und Messung an
den Sensoren ist das Innovationssignal.
\end{beispiel}

\section{Der strikte Modus}

Mit \texttt{RTWATERFLOW\_EXPOSE\_GROUND\_TRUTH=false} wird die Realitätsschicht
vollständig zurückgehalten: die Wahrheits-Schlüssel, Hintergrund-Leckagen und alle
Alarme werden aus den Frames \emph{und} den Aufzeichnungen entfernt — es bleibt
nur die Anlagen-SCADA, die ein realer Betreiber sähe. Die Schätzung überlebt den
strikten Modus, weil sie aus Messungen und nicht aus Wahrheit abgeleitet ist.

% ===================================================================
\chapter{Die Alarmzentrale (Compliance)}
\label{ch:alarm}

Nach jeder Rechnung prüft eine Regelpass die DVGW-Vorgaben und meldet typisierte
Befunde mit deutscher Zitatstelle (\file{docs/COMPLIANCE.md} bildet jede Regel
ab). Ein Befund trägt eine Schwere (\textcolor{rot}{\textbf{Verletzung}} /
\textcolor{amber}{\textbf{Warnung}} / Info) und einen Nachhaltigkeitszähler.

\begin{center}
\small
\begin{tabularx}{\linewidth}{@{}lX@{}}
\toprule
Prüfung & DVGW-Grundlage \\
\midrule
Mindestversorgungsdruck & W~400-1: $2{,}0 + 0{,}35\cdot(\text{Geschosse}-1)$~bar \\
Ruhedruck & 8~bar Warnung / 10~bar (PN~10) Verletzung \\
Geschwindigkeit & $>2{,}0$~m/s (momentan Warnung, dauerhaft Verletzung) \\
Stagnation/Selbstreinigung & Hygiene (Stundenmittel $<5$~mm/s) \\
Behälterreserve/-umschlag & Löschreserve, Umschlagzeit $<24$~h \\
Löschwasser & W~405: 48/96/192~m³/h nach Nutzung, $\geq1{,}5$~bar \\
Brunnenalterung & W~130: Ergiebigkeitsverlust $>10\,\%$ \\
Wasserrecht & WHG §§8--10 (Tages-Warnung / Jahres-Verletzung) \\
\bottomrule
\end{tabularx}
\end{center}

\begin{beispiel}[Keine Alarmflut]
Gesunde Netze lösen keine Flut aus: die per-Rohr-Hygienewarnungen eines fein
dimensionierten ländlichen Netzes bündeln sich oberhalb einer Schwelle zu
\emph{einem} Flotten-Befund (Selbstreinigung), sodass gesetzte Anomalien
hervorstechen. Die gemessene und die geschätzte Sicht erfinden nie ein falsches
„alles in Ordnung“. Die Compliance läuft in einem eigenen Fehlerkäfig — eine
gestörte Prüfung wird zum System-Info-Befund, verwirft aber nie den gerechneten
Takt.
\end{beispiel}

% ===================================================================
\chapter{Behälter, Pumpen und Druckzonen}
\label{ch:behaelter}

\section{Hochbehälter}
Ein Hochbehälter ist ein Festdruckknoten mit füllstandsintegrierendem Regler: die
Höhe wird vor der Rechnung als Druck geschrieben und nach der Rechnung explizit
(Euler) aus dem Zufluss fortgeschrieben (EPANET-Semantik), mit Überlauf- und
Leerlauf-Klemmen. Kennzahlen: nutzbares Volumen, Löschreserve (mit Alarm bei
Unterschreitung) und Pufferzeit.

\section{Pumpwerke}
Ein Pumpwerk folgt einer Q-H-Kennlinie mit Zweipunktregelung
(Ein-unter/Aus-über einem Behälter-Füllstand). Ein numerischer Kernpunkt: da
pandapipes den Kennlinienhub \emph{explizit je Newton-Schritt} anwendet und
Rückwärtsfluss widerstandsfrei durchlässt, drainiert eine laufende Pumpe gegen
dominante statische Höhe den Behälter rückwärts. Die Plattform zeigt dem Solver
deshalb einen \emph{konstanten} Hub und findet den ehrlichen Betriebspunkt über
eine eingeklammerte Sekante — mit Rückschlagorgan-Schließung, wenn eine Pumpe am
Nullförderpunkt umkehrt.

\begin{beispiel}[Operatormodi]
Jede Station kennt die Modi \ui{Auto} (Zweipunktregelung), \ui{Ein} und
\ui{Aus}. Der Modus wird szenariofest gespeichert; eine Rückschlag-Schließung
kann sich nie in \ui{Auto} festhaken, weil die Regel ihr Laufgedächtnis selbst
führt (statt es aus \texttt{in\_service} zurückzulesen).
\end{beispiel}

\section{Druckminderer und Druckzonen}
Ein Druckminderer (\texttt{press\_control}) hält am Auslauf einen Solldruck und
bildet die Grenze zwischen zwei \textbf{Druckzonen}. Er muss eine \emph{Schnitt}-Kante
sein (kein Umgehungsrohr) — sonst sähe die Tiefzone zwei widersprüchliche Höhen.
Ohne ihn läge eine tief gelegene Zone über 10~bar (PN~10); mit ihm bleibt sie im
Band (Abbildung~\ref{fig:drucklinie}).

\begin{figure}[t]
\centering
\begin{tikzpicture}[font=\small, node distance=9mm,
  n/.style={circle, draw=wasser, fill=wasserhell, minimum size=6mm, inner sep=0pt}]
\node[n,label=above:{HB 420 m}] (hb) {};
\node[n, right=14mm of hb] (h1) {};
\node[n, right=of h1] (h2) {};
\node[n, below right=6mm and 8mm of h2, fill=amber!30, draw=amber,
      label=below:{Druckminderer}] (prv) {};
\node[n, right=14mm of prv] (t1) {};
\node[n, right=of t1] (t2) {};
\begin{scope}[on background layer]
  \node[fit=(hb)(h1)(h2), draw=gruen, dashed, rounded corners, inner sep=4mm,
        label=above:{Hochzone 4--6 bar}] {};
  \node[fit=(t1)(t2), draw=wasser, dashed, rounded corners, inner sep=4mm,
        label=below:{Tiefzone}] {};
\end{scope}
\draw[thick] (hb)--(h1) (h1)--(h2) (h2)--(prv) (prv)--(t1) (t1)--(t2);
\end{tikzpicture}
\caption{Zwei Druckzonen, getrennt durch einen Druckminderer als Schnitt-Kante —
das Grundmuster von Musterdorf und (mehrfach) Neubeuern.}
\label{fig:druckzonen}
\end{figure}

% ===================================================================
\chapter{Nachfrage und Wetter}
\label{ch:nachfrage}

Die Nachfrage-Engine baut aus Archetyp und Größe realistische Verbrauchsprofile:
Tagesgang, Wochenendverhalten, Saison, Temperatur-Kopplung und die nächtliche
Freibad-Rückspülung (DIN~19643). Verbraucher ohne Größenangabe folgen einem
einfachen Ersatzprofil (\texttt{demand\_factor}).

\begin{center}
\small
\begin{tabularx}{\linewidth}{@{}lX@{}}
\toprule
Archetyp & Charakter \\
\midrule
\texttt{residential\_city} / \texttt{\_village} & Wohnen mit Morgen-/Abendspitze, Wochenendverschiebung \\
\texttt{industry} & Schichtblock \\
\texttt{school} & Vormittagspulse \\
\texttt{office} & Bürotag \\
\texttt{hospital} & gleichmäßig hoch \\
\texttt{farm\_dairy} / \texttt{\_pigs} & Melkpulse 05--07/16--18~Uhr, temperaturgekoppelt \\
\texttt{pool} & Freibad mit Mai--September-Saison und Rückspülung \\
\bottomrule
\end{tabularx}
\end{center}

\section{Der Hitzetag}
Der Umschalter \ui{Hitzetag} (Temperaturaufschlag + Trockenheit) verschiebt die
Spitze in den Abend und hebt die Nachtflüsse — der Trocken-Heiß-Effekt 2018 mit
abendlichem Bewässerungsblock. Die Umwelt-Sektion zeigt drei ehrliche Zustände
(\ui{Normal}, \ui{Hitzetag}, \emph{benutzerdefinierte Übersteuerung}) und
schützt gegen Wettläufe mit dem 5-Sekunden-Poll.

\section{Validierung gegen W~410}
Die synthetische Jahresstatistik liegt im DVGW-W-410-Korridor: Tages- und
Stundenspitzenfaktor $f_d = 3{,}9\cdot E^{-0{,}0752}$ bzw.\
$f_h = 18{,}1\cdot E^{-0{,}1682}$ innerhalb $\pm20\,\%$ (siehe
Anhang~\ref{app:validierung}).

% ===================================================================
\chapter{Ereignisse: Löschwasser, Rohrbruch, Leckage}
\label{ch:ereignisse}

Im Bereich \ui{Ereignisse} legen Sie druckabhängige Entnahmen an. Alle folgen
einem Emitter-Gesetz $\dot m = C\cdot\max(p,0)^{N_1}$; PDA und Emitter teilen sich
einen äußeren Fixpunkt.

\begin{description}[leftmargin=1.2em]
\item[Hydrant] ein W-405-Löschfall, dimensioniert auf eine Zielentnahme; auf dem
  Draht als Anlagen-Telemetrie sichtbar (auch strikt).
\item[Rohrbruch] Ausfluss über eine Bruchfläche ($C=C_d\,A\sqrt{2\rho}$, $C_d=0{,}75$).
\item[Leckage] eine FAVAD-Emitter ($N_1=1{,}15$); \emph{verborgene} Realität — die
  nächtliche Mindestwassermenge, die der Betreiber erst erkennen muss (im strikten
  Modus ausgeblendet).
\end{description}

\section{Druckabhängige Nachfrage (PDA)}
Der Umschalter \ui{PDA} macht den Kontrast „warum druckabhängige Nachfrage“
sichtbar. Nach Wagner liefert ein Anschluss bei $p\le p_{\min}$ nichts, bei
$p\ge p_{\text{req}}$ voll und dazwischen einen Anteil
$\bigl((p-p_{\min})/(p_{\text{req}}-p_{\min})\bigr)^{1/2}$. Unterversorgte
Anschlüsse liefern so \emph{weniger} Wasser statt unphysikalisch negativer Drücke;
ein Physik-Wächter stuft einen selbstkonsistenten Zustand mit negativem Druck
ehrlich auf \texttt{degraded} herab.

\begin{beispiel}[Ehrliche Verknappung]
Eine Großentnahme, die das Netz nicht decken kann, führt zu Teilbelieferung, einer
positiven Fehlmenge und einem stabilen Gleichgewicht bei rund 1{,}5~bar — kein
negativer Druck, kein Serverfehler. Die gelieferte Menge ist stets
Wagner-konsistent zum gemeldeten Druck.
\end{beispiel}

% ===================================================================
\chapter{Brunnen und Grundwasser}
\label{ch:brunnen}

Die \textbf{Rohwasserseite} ist reines Python: ein \emph{Reinwasserbehälter}
(Break Tank) entkoppelt sie hydraulisch vom pandapipes-Netz, sodass Brunnen +
Grundwasserleiter + Bilanzierung eine Massenbilanz sind, keine Hydraulikrechnung.

\begin{description}[leftmargin=1.2em]
\item[Grundwasserleiter] ein linearer Speicher: $\mathrm{d}h/\mathrm{d}t =
  (\text{Anreicherung}\cdot\text{Dürre} - \sum Q)/(S_y A)$, saisonale
  Anreicherung, Dürrefaktor.
\item[Brunnen] Absenkung $Q/(Q/s)$, Filterschutz (deckelt die Ergiebigkeit),
  Alterung (W~130, schneller bei tiefer Absenkung), Regeneration auf 95\,\%,
  Sichardt-Interferenz.
\item[Brunnenfeld] speist den Break Tank, bucht das Wasserrecht (WHG) und die
  Energie in kWh/m³.
\end{description}

\begin{figure}[t]
\centering
\begin{tikzpicture}[font=\small, node distance=7mm,
  b/.style={draw, rounded corners=2pt, minimum width=24mm, minimum height=8mm,
            align=center}]
\node[b, fill=wasserhell, draw=wasser] (aq) {Grundwasser-\\leiter};
\node[b, right=of aq, fill=wasserhell, draw=wasser] (well) {Brunnen};
\node[b, right=of well, fill=gray!8, draw=gray!60] (bt) {Reinwasser-\\behälter};
\node[b, right=of bt, fill=amber!20, draw=amber] (np) {Netzpumpe};
\node[b, right=of np, fill=wasserhell, draw=wasser] (hb) {Hochbehälter};
\node[b, right=of hb, fill=gruen!12, draw=gruen] (h) {Haushalte};
\foreach \a/\b in {aq/well, well/bt, bt/np, np/hb, hb/h}
  \draw[-{Latex[length=2mm]}, thick, wasser] (\a) -- (\b);
\end{tikzpicture}
\caption{Die Lauenau-Kette. Dürre $\to$ fallender Grundwasserstand $\to$
gedeckelte Brunnenergiebigkeit $\to$ leerer Break Tank $\to$ die Netzpumpe fällt
bei Trockenlauf aus $\to$ der Hochbehälter entleert sich $\to$ Haushalte fallen
trocken.}
\label{fig:kaskade}
\end{figure}

Der Dürre-Schieber und die Regenerier-Schaltfläche je Brunnen machen die Kaskade
(Abbildung~\ref{fig:kaskade}) im Zeitraffer erlebbar (Kapitel~\ref{ch:rundgaenge},
Lauenau).

% ===================================================================
\chapter{Der NetzStudio-Editor}
\label{ch:editor}

Der \ui{Editor} zeichnet ein Wassernetz auf \textbf{echten OpenStreetMap-Straßen}:
Quelle, Knoten und Verbraucher setzen (die EU-DEM-Höhe wird je Klick geholt),
Rohre straßengeführt ziehen (Einrasten + Dijkstra), aus dem Rohrkatalog wählen
(PE d110/d125/d160, GGG DN100/DN150/DN200).

\section{Live-Prüfung: die drei W-400-1-Lastfälle}
Vor der Übernahme prüft der Editor das gezeichnete Netz gegen die drei
Sizing-Lastfälle:

\begin{description}[leftmargin=1.2em]
\item[LF1 Maximale Förderung] Spitzendurchfluss, Pumpen an $\to$ Geschwindigkeit
  $\le2{,}0$~m/s.
\item[LF2 Spitzenstunde Maximaltag] jeder Anschluss $\ge$ seinem Geschossdruck
  ($2{,}0+0{,}35\cdot(\text{Geschosse}-1)$~bar), $\le8$~bar Ruhedruck.
\item[LF3 Löschfall] eine W-405-Löschentnahme am hydraulisch schlechtesten Knoten,
  Fließdruck $\ge1{,}5$~bar.
\end{description}

Erst ein \emph{bestehendes} Netz lässt sich in die Bibliothek übernehmen; die
Übernahme nutzt denselben Loader-Vertrag wie ein Import.

% ===================================================================
\chapter{Der Geodaten-Bündelbauer}
\label{ch:geodaten}

Das Offline-Werkzeug \file{tools/bundle\_builder} verwandelt eine echte deutsche
Stadt in ein Fünf-Datei-Bündel. Es ist zweistufig:

\begin{center}
\small
\begin{tabularx}{\linewidth}{@{}lX@{}}
\toprule
Schritt & Was passiert \\
\midrule
\texttt{make\_snapshot} (online) & OSM-Straßengraph holen (osmnx), nach UTM
  (EPSG:25832/25833) projizieren, Höhen per EU-DEM (OpenTopoData) oder lokalem
  DGM-GeoTIFF abtasten, alles in einen \emph{gepinnten} JSON-Schnappschuss frieren \\
\texttt{build\_from\_snapshot} (offline) & aus dem Schnappschuss deterministisch
  das Netz synthetisieren (nur networkx + Stdlib) \\
\bottomrule
\end{tabularx}
\end{center}

Der Offline-Schritt ist byte-stabil reproduzierbar — genau das ist das
M8-Akzeptanzkriterium. Aus dem Straßengraphen wird ein Schwerkraftnetz: die größte
zusammenhängende Komponente $\to$ ein Baum (Minimalspannbaum flach, ein
Schwerkraft-Abstiegsbaum hügelig) $\to$ eine \texttt{ext\_grid}-Quelle auf dem
höchsten Knoten $\to$ Verbraucher auf Blatt-/Straßenknoten $\to$ verjüngende
Durchmesser. Für ein hügeliges Netz fügt \texttt{enable\_prv\_zoning} Druckminderer
ein, wo der Talwärtsdruck die PN-10-Grenze überschritte. Details in
\file{docs/GEODATA\_BUILDER.md}.

% ===================================================================
\chapter{Szenarien-Rundgänge}
\label{ch:rundgaenge}

Jeder Rundgang nennt, was zu tun ist, \textbf{was man sieht} und \textbf{warum}.

\begin{beispiel}[Tutorial Hanglage — Hydrostatik]
Starten und einen Tag laufen lassen. \emph{Beobachtung:} Die Einspeisung
(„Hochbehälter“, 400~m) liegt bei 0{,}5~bar; am 54~m tiefer gelegenen Abnehmer
stehen $\sim5{,}78$~bar an. Der \textbf{Schlechtpunkt} ist nicht der tiefste,
sondern der \emph{höchste} Abnehmer ($\sim361$~m). \emph{Lehrpunkt:} $\sim1$~bar je
10~m Höhe — die Geländehöhe, nicht die Entfernung, bestimmt den Versorgungsdruck.
\end{beispiel}

\begin{beispiel}[Musterdorf — zwei Druckzonen]
Die Drucklinie zum tiefsten Abnehmer öffnen. \emph{Beobachtung:} Am Druckminderer
Talstraße fällt der Druck als sichtbare Stufe auf 2{,}8~bar; ohne ihn läge die
Tiefzone über 10~bar. Beide Zonen-Mediane liegen im 4--6-bar-Band. Der Hitzetag
verschiebt die Spitze auf $\sim$19:45. \emph{Lehrpunkt:} Druckzonen schützen die
tief gelegenen Rohre vor Überdruck.
\end{beispiel}

\begin{beispiel}[Mustertal — Gegenbehälter]
Den Abschnitt zum Wasserturm über den Tag beobachten. \emph{Beobachtung:} Der
Durchfluss kehrt die Richtung um ($\sim+5{,}7$~kg/s tags, $\sim-3{,}2$~kg/s in der
Schwachlast), während das Schaltband 1{,}2/3{,}4~m gehalten wird. \emph{Lehrpunkt:}
Ein Gegenbehälter am Netzende füllt sich nachts und stützt tagsüber die Spitze.
\end{beispiel}

\begin{beispiel}[Lauenau — Dürrekaskade]
Den Dürre-Schieber hochziehen und mehrere Tage schnell durchlaufen.
\emph{Beobachtung:} Der Grundwasserstand fällt, die Brunnenkapazität sinkt unter
die Spitze, der Reinwasserbehälter leert sich, die Netzpumpe fällt bei Trockenlauf
aus, der Hochbehälter entleert sich — die Haushalte fallen trocken (PDA). Am
Normaltag: 0{,}385~kWh/m³. \emph{Lehrpunkt:} die Kette Grundwasser $\to$ Brunnen
$\to$ Behälter $\to$ Netz und wo sie reißt (Abbildung~\ref{fig:kaskade}).
\end{beispiel}

\begin{beispiel}[Alpen — reales flaches Netz]
\emph{Beobachtung:} 182 echte Knoten, Drücke über den Tag $\sim2{,}8$--$5{,}9$~bar,
Spitzengeschwindigkeit $\sim0{,}87$~m/s in der Morgenspitze (weit unter der
2{,}0-m/s-Warnung), keine Verletzungen; auf der Karte der OSM-/EU-DEM-Quellen\-nachweis.
\emph{Lehrpunkt:} aus echten Straßen + DGM synthetisierte, solide Netze.
\end{beispiel}

\begin{beispiel}[Neubeuern — hügelige Druckzonen]
\emph{Beobachtung:} 215 echte Knoten über $\sim78$~m Relief, sechs Druckminderer
teilen die Oberstadt von der Talsohle; Drücke $\sim2{,}5$--$6{,}8$~bar, keine Flut.
Der Löschfall-Lastcheck (LF3) \emph{scheitert} am Schlechtpunkt. \emph{Lehrpunkt:}
echtes Relief erzwingt echte Druckzonen — und ein Schwerkraftnetz liefert nicht
überall die Löschwassermenge.
\end{beispiel}

\begin{beispiel}[Kevelaer — Stadtmaßstab]
\emph{Beobachtung:} $\sim540$ Knoten, in-band (4{,}3--5{,}2~bar), warmer Takt
$\sim37$~ms (Beobachter aus). \emph{Lehrpunkt:} die Rechenzeit skaliert mit der
Zahl der Solver-Aufrufe, nicht mit der Knotenzahl.
\end{beispiel}

% ===================================================================
\chapter{Aufzeichnung, Export und Szenarien}
\label{ch:export}

\section{Szenarien speichern und laden}
Ein Szenario-Rezept fasst Netz, Betriebshandlungen, Messgeräte, Hydraulik-Einstellungen
(PDA, Leckagen, Hydranten) und die Uhr in einer bearbeitbaren JSON-Datei zusammen.
Beim Laden werden die Einträge tolerant wiedergegeben — fehlt ein Element, wird der
Rest angewandt.

\section{Aufzeichnung und Export}
\ui{Aufzeichnung} schreibt jeden veröffentlichten Schritt als
CSV-Paket (eine Datei je Wire-Tabelle) plus ein Metadaten-Rezept. \ui{Tage
exportieren} spielt die aktuelle Konfiguration offline durch — \textbf{byte-kompatibel}
zur Live-Aufzeichnung: der Export startet seine gleitenden Fenster und die
Rohwasserseite ab Mitternacht neu und schaltet den (nicht-deterministischen)
Beobachter ab, sodass ein exportierter Tag eine Live-Aufzeichnung desselben Tages
datei-genau reproduziert.

\begin{beispiel}[CSV-Paket auswerten]
Die Pakete lassen sich direkt mit pandas lesen:
\begin{lstlisting}[language=Python]
import pandas as pd
consumers = pd.read_csv("data/recordings/<sitzung>/consumers.csv")
tanks     = pd.read_csv("data/recordings/<sitzung>/tanks.csv")
\end{lstlisting}
\end{beispiel}

% ===================================================================
\chapter{Langzeit-Visualisierung mit Grafana}
\label{ch:grafana}

Im Docker-Stack schreibt der \emph{collector}-Dienst die Telemetrie im
Line-Protocol nach InfluxDB (Port~8088); Grafana (Port~3002) zeigt sie als
Zeitreihen-Dashboards. So lassen sich Behälterstände, Zonendrücke und
Brunnenergiebigkeit über viele simulierte Tage verfolgen — jenseits des
gleitenden Verlaufspuffers der Oberfläche.

\begin{beispiel}[Zugangsdaten]
Das Grafana-Dashboard ist datei-provisioniert und läuft mit Entwickler-Zugangsdaten
gegen ein unauthentifiziertes Backend — bewusst für den lokalen Gebrauch. Vor einer
Freigabe ins Netz unbedingt Zugangsdaten ändern und absichern.
\end{beispiel}

% ===================================================================
\chapter{Konfiguration}
\label{ch:konfiguration}

Alle Einstellungen sind Umgebungsvariablen mit dem Präfix \texttt{RTWATERFLOW\_}:

\begin{center}
\small
\begin{tabularx}{\linewidth}{@{}llX@{}}
\toprule
Variable & Standard & Bedeutung \\
\midrule
\texttt{PORT} & 8002 & Backend-Port (Schwester-Schema) \\
\texttt{DEFAULT\_NETWORK} & tutorial\_hillside & Startnetz \\
\texttt{STEP\_INTERVAL\_SECONDS} & 1{,}0 & Wandzeit je Simulationsminute \\
\texttt{STEPS\_PER\_DAY} & 1440 & Schritte je Simulationstag \\
\texttt{AUTOSTART} & true & Engine beim Start laufen lassen \\
\texttt{EXPOSE\_GROUND\_TRUTH} & true & false = strikter Modus \\
\texttt{PDA\_ENABLED} & true & druckabhängige Nachfrage (false = Festentnahme-Kontrast) \\
\texttt{SOLVER\_ITER} & 100 & Basis-Iterationen der Solver-Leiter \\
\texttt{RECORD} & false & Daueraufzeichnung \\
\bottomrule
\end{tabularx}
\end{center}

% ===================================================================
\chapter{Grenzen des Modells}
\label{ch:grenzen}

\begin{description}[leftmargin=1.2em]
\item[Quasistatisch] Jeder Schritt ist eine stationäre Hydraulikrechnung;
  Druckstöße (DVGW~W~303) werden nicht berechnet.
\item[Druckminderer ohne Zustandsautomat] \texttt{press\_control} hält einen
  Sollwert, hat aber keine echte Ventilmechanik; Zonengrenzen sind statische
  Schnitte.
\item[Keine Wassergüte] Wasseralter/Stagnation erscheinen nur als Hygiene-Befund,
  nicht als Stofftransport.
\item[Beobachter nicht deterministisch] die Schätzung nutzt eine Wandzeit-Drosselung
  und ist beim Offline-Export abgeschaltet, damit der Export byte-stabil bleibt.
\item[Keine Feldvalidierung] die Lehrnetze sind synthetisch bzw.\ aus Geodaten
  synthetisiert — validiert ist die \emph{Engine} (Anhang~\ref{app:validierung}),
  nicht ein konkretes reales Versorgungsgebiet.
\end{description}

Die frühen M0-Grenzen (feste Entnahmen, keine Behälterdynamik) sind längst
aufgehoben: PDA, Behälter/Pumpen, Brunnen und der Beobachter sind aktiv.

% ===================================================================
\appendix
\chapter{REST-/WebSocket-Schnittstelle}
\label{app:api}

Das Backend bietet 65~Routen; die vollständige Referenz steht in
\file{docs/API.md} und interaktiv unter \api{/docs} (Swagger). Ein Auszug:

\begin{center}
\small
\begin{tabularx}{\linewidth}{@{}llX@{}}
\toprule
Methode & Route & Zweck \\
\midrule
GET  & \api{/state}      & der aktuelle Frame (gepollt) \\
GET  & \api{/history}    & der gleitende Verlauf \\
WS   & \api{/ws}         & der Live-Frame-Strom \\
GET  & \api{/network}    & Topologie + Anschlusswerte \\
GET  & \api{/findings}   & Alarme + Schwere-Zählung \\
GET  & \api{/tanks} / \api{/stations} & Behälter- / Pumpen-SCADA \\
POST & \api{/station/\{name\}} & Pumpenmodus auto/on/off \\
GET  & \api{/wellfields} & Brunnenfeld-SCADA \\
POST & \api{/wellfield/drought} & Dürrefaktor setzen \\
POST & \api{/pda}        & PDA an/aus \\
POST & \api{/hydrant} / \api{/burst} / \api{/leakage} & Ereignis anlegen \\
GET  / POST & \api{/environment} & Wetter lesen/setzen (Hitzetag) \\
POST & \api{/networks/import} & eigenes Bündel übernehmen \\
GET  & \api{/manual}     & dieses Handbuch (Kurzfassung) \\
\bottomrule
\end{tabularx}
\end{center}

\begin{lstlisting}[language=bash]
# aktuellen Zustand abfragen
curl http://localhost:8002/api/state
# eine Pumpe von Hand einschalten
curl -X POST http://localhost:8002/api/station/Pumpwerk -d '{"mode":"on"}'
# den Live-Strom mitlesen
websocat ws://localhost:8002/ws
\end{lstlisting}

\chapter{Eingabedateiformat}
\label{app:dateiformat}

Fünf JSON-Dokumente je Netz, beim Laden pydantic-geprüft (Knotenreferenzen,
Ganztags-Horizont, mindestens eine Druckquelle je zusammenhängendem Netz — ein
ext\_grid oder Behälter; mehrere sind seit M2 zulässig —, Erreichbarkeit aller
Knoten, keine isolierten Knoten, Druckminderer und Pumpen als Schnitt-Kanten).
Ein minimales Bündel:

\begin{lstlisting}[basicstyle=\ttfamily\scriptsize]
// network_structure.json
{"name":"Mini","junctions":[
  {"name":"hb","kind":"source","elevation_m":100,"pn_bar":0.5,"geo":[51.5,6.5]},
  {"name":"a", "kind":"consumer","elevation_m":60,"pn_bar":4.0,"geo":[51.49,6.51]}]}
// pipes.json
{"pipes":[{"from_node":"hb","to_node":"a","dn":110,"material":"PE",
           "geometry":[[51.5,6.5],[51.49,6.51]]}]}
// consumers.json
{"consumers":[{"node":"a","name":"Dorf","mdot_kg_per_s":0.2,
               "kind":"residential_village","storeys":2,"size":{"population":80}}]}
// supply.json
{"supplies":[{"node":"hb","name":"Einspeisung","kind":"ext_grid","p_bar":0.5}],
 "prvs":[],"tanks":[],"stations":[],"wellfields":[]}
// environment.json
{"resolution_minutes":15,"steps":96,"t_air_c":[...],"day_types":["workday"],
 "dryness":[0.3],"season_day_of_year":205}
\end{lstlisting}

\chapter{Validierung im Detail}
\label{app:validierung}

Die Hydraulik-Engine ist gegen das unabhängige Referenzwerkzeug EPANET (über
WNTR) kreuzvalidiert. Die Ergebnisse (voller Text in \file{docs/BENCHMARKS.md}):

\begin{center}
\small
\begin{tabularx}{\linewidth}{@{}lXl@{}}
\toprule
Prüfung & Referenz & Ergebnis \\
\midrule
Netzhydraulik & EPANET Net1 (9~Knoten) + Net3 (92) mit \api{swamee-jain}
  & Druck $<0{,}001$~bar, Durchfluss $\sim0{,}5\,\%$ \\
Höhendruck & pandapipes-Höhenbeispiel & 0{,}5~bar@400~m $\to$ 5{,}78~bar; $\pm0{,}01$~bar \\
Hydrant & EPANET/WNTR-Emitter & Knotendruck $\pm0{,}1$~bar \\
Behälter & EPANET/WNTR-Behälter & Füllraten-Median $<1$~cm/Takt \\
Nachfrage (PDA) & WNTR-PDD & Wagner-Kurve auf 0{,}001 \\
Nachfrage-Hüllkurve & DVGW W~410 $f_d/f_h$ & innerhalb $\pm20\,\%$ \\
Pumpenergie & TF §5 Korridor & Lauenau $\approx0{,}385$~kWh/m³ \\
\bottomrule
\end{tabularx}
\end{center}

Die Standardnetze werden dabei als Schwerkraftnetze nachgebaut (jeder Behälter/jede
Quelle als Festhöhe, Pumpen weggelassen — pandapipes modelliert eine Pumpe als
Konstant-Hub außerhalb des Newton-Solves; die Pumpen-/Behälter-\emph{Regelung} ist
über das Behälter-Orakel getrennt geprüft). Der Druck wird mit pandapipes'
eigenem statischem Gradienten in bar umgerechnet, sodass die Gegenüberstellung die
hydraulische Höhe vergleicht und nicht einen $\sim0{,}1\,\%$-Dichte-Versatz.

Reproduktion (im \texttt{[dev]}-Venv, das \texttt{wntr} enthält):
\begin{lstlisting}[language=bash]
python -m pytest tests/validation -q                     # EPANET Net1/Net3 + PDA
python -m pytest tests/test_tank_oracle.py \
                 tests/test_hydrant_oracle.py -q          # Behälter/Hydrant
\end{lstlisting}

\chapter{Fehlerbehebung}
\label{app:fehler}

\begin{description}[style=nextline,leftmargin=1.2em]
\item[Der Beobachter läuft, aber die Schätzung wirkt „eingefroren“.]
  Der Beobachter drosselt sich über die Wandzeit selbst und ist beim Export
  bewusst abgeschaltet — im Live-Betrieb aktualisiert er mit eigenem Raster.
\item[Ein neuer Wasserzähler zeigt nichts an.]
  Ehrlicher Kaltstart: im 15-Minuten-Modus liefert ein Zähler erst nach dem ersten
  Fensterabschluss einen Mittelwert.
\item[Ein Takt meldet \texttt{degraded}.]
  Die Solver-Leiter musste auf eine gröbere Stufe (swamee-jain/nikuradse)
  ausweichen, oder der Physik-Wächter hat einen unphysikalischen Zustand
  herabgestuft — das ist Datum, kein Absturz; der nächste warm gestartete Takt
  fängt sich meist.
\item[Docker: das Backend startet nicht.]
  Prüfen Sie mit \texttt{docker compose logs backend}; Ports 8002/8082/8088/3002
  dürfen nicht belegt sein.
\end{description}

\end{document}
